% Combined UK Patent Application — Specification
% Application 1: AMTTP + Adaptive Friction
% Applicant: Odeyemi Olusegun Israel
% Date: 4 March 2026

\documentclass[12pt,a4paper]{article}
\usepackage[margin=2.5cm]{geometry}
\usepackage{amsmath,amssymb}
\usepackage{graphicx}
\usepackage{enumitem}
\usepackage{hyperref}
\usepackage{titlesec}
\usepackage{fancyhdr}
\usepackage{booktabs}
\setlength{\headheight}{14.5pt}

\pagestyle{fancy}
\fancyhf{}
\rhead{Patent Application}
\lhead{AMTTP/AF --- Odeyemi O.I.}
\rfoot{Page \thepage}

\titleformat{\section}{\large\bfseries}{\thesection.}{0.5em}{}
\titleformat{\subsection}{\normalsize\bfseries}{\thesubsection.}{0.5em}{}
\titleformat{\subsubsection}{\normalsize\itshape}{\thesubsubsection.}{0.5em}{}

\begin{document}

\begin{center}
{\LARGE\bfseries PATENT SPECIFICATION}\\[1em]
{\Large State-Adaptive Compliance Enforcement System\\
for Decentralised Finance Networks}\\[2em]
{\large Applicant: Odeyemi Olusegun Israel}\\[0.5em]
{\large Inventor: Odeyemi Olusegun Israel}\\[0.5em]
{\large Date of Filing: 4 March 2026}\\[0.5em]
{\large Application Number: [TO BE ASSIGNED]}\\[0.3em]
{\normalsize \textit{Filed pursuant to the Patents Act 1977}}
\end{center}

\vspace{2em}

% ============================================================
\section{TECHNICAL FIELD}
% ============================================================

The present invention relates to systems and methods for enforcing
anti-money laundering (AML) and regulatory compliance in decentralised
finance (DeFi) transactions using blockchain smart contracts, machine
learning risk scoring, and zero-knowledge cryptographic proofs, and more
particularly to such systems in which enforcement parameters are computed
in real time from a closed-form state-adaptive controller that models the
transaction processing network as a multi-body dynamical system with an
associated energy landscape.

% ============================================================
\section{BACKGROUND}
% ============================================================

\subsection{Technical Problems}

Decentralised finance protocols process financial transactions on blockchain
networks without centralised intermediaries. This architecture creates two
interacting technical problems that prior art has not solved in combination.

\subsubsection{Pre-Settlement Compliance Enforcement}

Existing blockchain analytics platforms (Chainalysis, Elliptic) operate
retrospectively: they identify suspicious transactions \emph{after}
settlement, by which time irrecoverable value transfer has already occurred.
Platform-level KYC/AML checks are bypassed entirely when users transact
directly through a DeFi protocol. Supervised fraud detection models require
large quantities of labelled fraudulent transactions; in practice fewer than
0.1\% of blockchain transactions are labelled as illicit, producing severe
class imbalance. Privacy regulations (UK Data Protection Act 2018, EU GDPR)
conflict with the traditionally disclosure-intensive compliance model.

\subsubsection{Static Enforcement Parameters}

Existing compliance systems apply fixed risk thresholds, fixed escrow rules,
and fixed proof requirements regardless of the structural state of the
underlying transaction network. This is provably suboptimal: above a critical
manifold in state space, any constant enforcement configuration necessarily
fails to achieve a state-dependent welfare optimum. No prior art derives a
closed-form optimal state-adaptive enforcement function as a mapping from a
computed energy landscape to numerical thresholds and proof requirements.

\subsection{Prior Art Limitations}

Prior art known to the applicant includes blockchain analytics retrospective
monitoring platforms; exchange-level KYC/AML systems that do not extend to
protocol-level transactions; standard supervised fraud detection systems
requiring labelled datasets; zero-knowledge privacy protocols (e.g., Zcash)
that conceal transactions without implementing compliance; macroprudential
risk measures including CoVaR, SRISK, and the Basel III countercyclical
capital buffer rule (which applies a constant-friction piecewise-linear
function to the credit-to-GDP gap and does not derive a closed-form
state-dependent optimal policy); and agent-based financial network
simulations lacking closed-form optimal policies and provable stability
guarantees. None of the prior art provides the combination of pre-settlement
enforcement, weak supervision, closed-form state-dependent threshold and
proof-requirement selection, and privacy-preserving verification.

% ============================================================
\section{SUMMARY OF THE INVENTION}
% ============================================================

The invention provides a state-adaptive compliance enforcement system and
method for DeFi networks whose enforcement configuration is computed in
real time by a closed-form controller derived from a dynamical model of the
transaction network. The shared inventive concept is the
\emph{state-dependent selection of a risk threshold profile and a
proof-requirement policy as a function of network instability metrics
computed from an energy landscape derived by modelling the network as a
multi-body dynamical system}. The primary technical contributions are:

\begin{enumerate}[label=(\roman*)]
\item A \textbf{four-layer DeFi compliance architecture} comprising interface,
compliance orchestration, offline training, and blockchain infrastructure
layers that together implement a deterministic compliance lifecycle enforcing
policy actions \emph{before} transaction settlement.

\item A \textbf{weak supervision pipeline} comprising a composite teacher
ensemble (autoencoder--gradient boosting, FATF rule-based, graph-structural)
generating pseudo-labels from unlabelled blockchain data, and a student model
trained on those pseudo-labels, eliminating the need for labelled fraud data.

\item A \textbf{state-adaptive transaction flow control module} that computes
one or more network instability metrics from the transaction graph and
aggregate risk dynamics, and selects an enforcement profile from a
predetermined set, the selected profile deterministically defining:
\begin{enumerate}[label=(\alph*)]
\item a \emph{risk threshold profile} comprising numerical thresholds for the
compliance decision matrix; and
\item a \emph{proof-requirement policy} specifying, for at least one risk
score tier or transaction class, one or more cryptographic proofs to be
verified on-chain as a prerequisite to settlement.
\end{enumerate}

\item A \textbf{gravitational interaction engine} that models financial
institutions as particles in a feature space under a pairwise potential
comprising erf-attraction and logarithmic repulsion, with a blind-spot energy
coupling, producing a combined systemic energy function $E(X)$.

\item A \textbf{spectral phase transition detector} computing the spectral
radius of the interaction matrix and determining whether the network state
lies above or below a critical manifold, above which state-adaptive control
is provably necessary.

\item A \textbf{closed-form optimal adaptive policy} $\gamma^*(X)$ derived
from the Bellman equation, mapping systemic energy and its gradient to an
optimal friction coefficient, which is further mapped to the enforcement
profile defining the risk threshold profile and proof-requirement policy.

\item A \textbf{zero-knowledge proof subsystem} (zkNAF) comprising Groth16
zk-SNARK circuits for sanctions non-membership, risk score range, and KYC
verification proofs, enabling privacy-preserving on-chain verification.

\item \textbf{Smart contract infrastructure} that enforces all policy actions
deterministically before settlement, with an immutable on-chain audit trail.
\end{enumerate}

% ============================================================
\section{DETAILED DESCRIPTION}
% ============================================================

\subsection{System Architecture Overview}

The system comprises four hierarchically organised layers:

\textbf{Layer~I --- Interface Layer.}
Client SDKs (TypeScript and Python) and a RESTful API provide the entry
points through which external applications submit transaction requests,
query risk scores, and verify compliance status. A web dashboard provides
visual compliance monitoring.

\textbf{Layer~II --- Compliance Orchestration Layer.}
The core compliance logic resides here. It comprises:
\begin{itemize}
\item A \emph{machine learning risk scoring engine} computing a numerical
risk score from 167 engineered features (address-level, transaction-level,
graph-level) for each submitted transaction request.
\item A \emph{graph analytics module} constructing transaction graphs and
computing centrality, clustering, PageRank and community features.
\item A \emph{sanctions screening module} checking counterparties against
OFAC and EU sanctions databases.
\item A \emph{state-adaptive transaction flow control module} (the adaptive
friction controller, described in \S\ref{sec:af}) computing instability
metrics and selecting the enforcement profile.
\item A \emph{deterministic policy engine} mapping calibrated risk scores to
policy actions via a decision matrix with thresholds drawn from the enforced
profile.
\end{itemize}

\textbf{Layer~III --- Offline Training Layer.}
Implements the composite teacher ensemble, pseudo-label generation, student
model training, cross-validation, and sigmoid score calibration.

\textbf{Layer~IV --- Infrastructure Layer.}
Blockchain smart contracts, the zkNAF proof subsystem, databases, and
distributed file storage for immutable audit records.

\subsection{Composite Teacher Weak Supervision Pipeline}

The offline training layer constructs three independent teacher models to
generate pseudo-labels from unlabelled transaction data:

\begin{enumerate}[label=(\alph*)]
\item \textbf{Autoencoder--Gradient Boosting Teacher ($s_{\text{AE-XGB}}$):}
An autoencoder learns a compressed representation of transaction features;
a gradient boosting model is trained on reconstruction errors and latent
features to estimate anomaly scores.

\item \textbf{FATF Rule-Based Teacher ($s_{\text{FATF}}$):} Encodes
Financial Action Task Force red-flag indicators (rapid fund movement,
structuring patterns, round-number transactions, high-risk jurisdiction
involvement) as deterministic scoring rules.

\item \textbf{Graph-Based Teacher ($s_{\text{Graph}}$):} Computes
suspicion scores from graph-structural properties: proximity to known
high-risk addresses, temporal clustering, and unusual connectivity.
\end{enumerate}

The composite teacher score is:
\begin{equation}
s_{\text{teacher}}(x) = w_1 s_{\text{AE-XGB}}(x) +
w_2 s_{\text{FATF}}(x) + w_3 s_{\text{Graph}}(x)
\end{equation}

In the preferred embodiment $w_1{=}0.4$, $w_2{=}0.3$, $w_3{=}0.3$.
Binary pseudo-labels are produced by thresholding and used to train an
XGBoost student model on 167 engineered features.

\subsection{Risk Score Calibration}

The raw student model output is recalibrated:
\begin{equation}
\hat{p}_{\text{cal}}(x) = \sigma\!\left(k\cdot(s(x)-P_{75})\right)
\end{equation}
where $\sigma$ is the logistic sigmoid, $k{=}30$ is a steepness parameter,
and $P_{75}$ is the 75th percentile of the training score distribution. The
calibrated probability is mapped to an integer score on $[0,1000]$.

\subsection{Compliance Decision Matrix}

The deterministic decision matrix maps integer scores to policy actions:

\begin{center}
\begin{tabular}{lll}
\toprule
\textbf{Score Range} & \textbf{Policy Action} & \textbf{Smart
Contract Effect} \\
\midrule
$[0,\,\tau_1)$ & APPROVE & Immediate execution \\
$[\tau_1,\,\tau_2)$ & APPROVE + MONITOR & Execution; monitoring flag \\
$[\tau_2,\,\tau_3)$ & REVIEW & Queue for manual review \\
$[\tau_3,\,\tau_4)$ & ESCROW & Time-locked escrow \\
$[\tau_4,\,1000]$ & BLOCK & On-chain rejection \\
\bottomrule
\end{tabular}
\end{center}

Under the \emph{stable} enforcement profile the thresholds are $(\tau_1,
\tau_2, \tau_3, \tau_4) = (200,400,600,800)$. The state-adaptive module
selects profiles that adjust these thresholds and impose additional proof
requirements as a function of computed network instability.

\subsection{State-Adaptive Transaction Flow Control Module}
\label{sec:af}

\subsubsection{Instability Metrics}

The module computes one or more network instability metrics from:
\begin{enumerate}[label=(\alph*)]
\item the spectral radius or centrality concentration of the transaction
graph constructed from recent on-chain activity over a sliding time window;
\item a temporal burst indicator computed from the transaction count process;
and/or
\item a shift of the calibrated risk score distribution relative to a
reference distribution.
\end{enumerate}

\subsubsection{Enforcement Profile Selection}

Based on the computed metrics the module selects one profile from a discrete
set (e.g., \textit{stable}, \textit{stressed}, \textit{critical}). Each
profile deterministically defines:
\begin{enumerate}[label=(\alph*)]
\item a \emph{risk threshold profile} $(\tau_1,\tau_2,\tau_3,\tau_4)$ for
the decision matrix; and
\item a \emph{proof-requirement policy} specifying, for at least one risk
score tier or transaction class, which of the zkNAF proofs (sanctions
non-membership, risk range, KYC) must be verified on-chain prior to
settlement.
\end{enumerate}

For a fixed system state the selected profile and all resulting decisions are
deterministic and reproducible.

\subsubsection{Macro-Level Systemic Energy Computation}

When financial metric data for multiple institutions is available (e.g.,
aggregated on-chain liquidity metrics, cross-protocol risk exposures), the
module further computes network instability using a gravitational interaction
model. The system models $N$ institutions as particles in $\mathbb{R}^d$
governed by:
\begin{equation}
\dot{X} = -\alpha(X-\mu\mathbf{1}^\top) - \nabla_X\Phi_{\text{pair}}(X)
- \sum_k \beta_k D\Phi_k(X)^\top D_k(X)
\end{equation}

where $X \in \mathbb{R}^{N\times d}$ is the state matrix, $\alpha$ is a
mean-reversion rate, $\Phi_{\text{pair}}$ is the pairwise interaction
potential, and $\Phi_k$ are external field operators.

\subsubsection{Pairwise Interaction Potential}

The pairwise potential between institutions $i$ and $j$ at distance
$r = \|x_i - x_j\|$ is:
\begin{equation}
\phi(r) = \gamma \cdot \frac{\sigma\sqrt{\pi}}{2}\,
\mathrm{erf}\!\left(\frac{r}{\sigma}\right)
- \gamma\lambda_{\mathrm{rep}} \cdot \log(r+\varepsilon)
\end{equation}

The erf-attraction is bounded (unlike gravitational $1/r$ attraction),
and the logarithmic repulsion prevents institution collapse. In the preferred
embodiment: $\alpha{=}0.10$, $\gamma{=}1.00$, $\sigma{=}1.00$,
$\lambda_{\mathrm{rep}}{=}0.10$, $\varepsilon{=}10^{-5}$.

\subsubsection{Blind-Spot Energy Module}

A complementary energy function is computed from four deviation channels
relative to reference distributions fitted on a normal period:
\begin{enumerate}[label=(\alph*)]
\item enstrophy anomaly $\delta_C$ --- rotational energy deviation;
\item spectral gap $\delta_G$ --- structural fragility indicator;
\item alignment anomaly $\delta_A$ --- co-movement deviation; and
\item temporal novelty $\delta_T$ --- drift from reference state patterns.
\end{enumerate}

The blind-spot energy is:
\begin{equation}
E_{\mathrm{BS}}(X) =
\|(X-\mu_0)^\top\Sigma_0^{-1}(X-\mu_0)\|_F
\end{equation}
where $\mu_0$ and $\Sigma_0$ are the reference mean and covariance estimated
by Ledoit--Wolf Oracle shrinkage.

\subsubsection{Critical Manifold}

A critical manifold $\mathcal{C}$ separates regimes in which constant
enforcement parameters are adequate from regimes in which they are provably
suboptimal:

\textbf{Theorem (Spectral Phase Transition):}
\begin{equation}
\mathcal{C} = \{X : \lambda_{\max}(\rho(X)\cdot\ell(X)\cdot W) = 1\}
\end{equation}

where $\rho(X)$ is the spectral radius of the inter-institution correlation
matrix, $\ell(X)$ is the system leverage factor, and $W$ is the weighted
adjacency matrix. Above $\mathcal{C}$ state-adaptive friction is both
necessary and sufficient for optimality.

\subsubsection{Closed-Form Optimal Adaptive Policy}

\textbf{Theorem (Optimal Friction):} The optimal state-dependent friction
coefficient solving the associated Bellman equation is:

\begin{equation}
\gamma^*(X) \approx
\frac{\beta\|\nabla E\|^2}{2\kappa + \beta\|\nabla E\|^2}\cdot
\frac{E(X)}{E(X)+\theta}
\end{equation}

where $E(X)$ is the combined systemic energy, $\kappa$ is the economic cost
parameter, $\beta$ is the systemic damage coefficient, and
$\theta=\kappa/\beta$. The policy satisfies $\gamma^*\to 0$ as
$E\to 0$ (no friction when stable) and $\gamma^*\to 1$ as $E\to\infty$ (maximum friction under severe instability).

The scalar $\gamma^*(X)\in[0,1]$ is used to select the enforcement profile:
below a lower threshold $\gamma_1$ the \textit{stable} profile is active;
between $\gamma_1$ and $\gamma_2$ the \textit{stressed} profile is active;
above $\gamma_2$ the \textit{critical} profile is active.

\subsubsection{Navier--Stokes Analogy (Enhanced Embodiment)}

Institution position changes $\dot{X}$ are treated as a velocity field.
The velocity gradient tensor is decomposed into strain $S$ and vorticity
$\Omega$. Enstrophy $\mathcal{E} = \frac{1}{2}\|\omega\|_F^2$ quantifies
rotational energy. The Beale--Kato--Majda criterion adapted to financial
networks serves as a singularity indicator. The adaptive friction coefficient
is reinterpreted as an adaptive viscosity:
$\nu = \nu_0(1+\gamma^*(E_{\mathrm{BS}}))$. Minsky regime classification
(hedge, speculative, Ponzi) emerges from the cosine alignment between the
gravitational gradient and the blind-spot gradient.

\subsection{Zero-Knowledge Proof Subsystem (zkNAF)}

The zkNAF subsystem implements three Groth16 zk-SNARK proof circuits:

\textbf{Sanctions Non-Membership Proof.} The prover demonstrates that a
blockchain address is not a member of a Merkle tree of sanctioned addresses
(using the Poseidon hash function) without revealing the address. The
verifier checks the proof against the published Merkle root.

\textbf{Risk Range Proof.} The prover demonstrates:
$s_{\min} \leq s(x) \leq s_{\max}$ without revealing the exact score $s(x)$.

\textbf{KYC Verification Proof.} The prover demonstrates completion of
Know Your Customer verification with an accredited provider without
revealing identity data.

Each proof has a 24-hour time-to-live. Verification is performed on-chain
by a dedicated ZK verifier smart contract that performs the bilinear pairing
check and records the result with a timestamp.

\subsection{Smart Contract Architecture}

\begin{itemize}
\item \textbf{Core Contract:} Implements the deterministic decision matrix
on-chain; registers compliance decisions as an immutable audit record.
\item \textbf{Escrow Contract:} Holds funds in time-locked escrow for
ESCROW-policy transactions; supports release and dispute escalation.
\item \textbf{ZK Verifier Contract:} Verifies Groth16 proofs; stores
binary verification results and timestamps.
\item \textbf{Cross-Chain Bridge Contract:} Requires compliance verification
on both source and destination chains before releasing bridged funds.
\item \textbf{Dispute Resolution Contract:} Structured escalation and
release process for escrowed transactions under dispute.
\end{itemize}

All enforcement decisions execute \emph{before} the corresponding
transaction is settled on the blockchain network.

\subsection{Feature Engineering}

The risk scoring engine operates on 167 features across three categories:
\begin{itemize}
\item \textbf{Address-level:} Transaction count, total value, mean and
variance of amounts, address age, unique counterparty count, temporal
activity patterns.
\item \textbf{Transaction-level:} Value, gas parameters, input data
characteristics, token transfer indicators.
\item \textbf{Graph-level:} Degree centrality, betweenness centrality,
clustering coefficient, PageRank, community membership, local graph density.
\end{itemize}

\subsection{Performance Characteristics}

\textbf{DeFi Compliance System (trained on 2.64 million Ethereum
transactions):}
\begin{center}
\begin{tabular}{ll}
\toprule
\textbf{Metric} & \textbf{Value} \\
\midrule
Meta-Ensemble ROC-AUC (test) & 1.0000 \\
Meta-Ensemble ROC-AUC (5-fold CV) & $0.9965\pm0.0049$ \\
Meta-Ensemble PR-AUC & 0.9999 \\
Transaction-Level XGBoost AUC & 0.9878 \quad F1: 0.901 \\
Cross-domain mean AUC (9 datasets) & 0.9962 \\
Escrow contract gas cost & 287,890 gas \\
\bottomrule
\end{tabular}
\end{center}

\textbf{Adaptive Friction Controller (FDIC SDI data, 140 quarters):}
\begin{center}
\begin{tabular}{lccc}
\toprule
\textbf{Configuration} & $N$ & \textbf{AUROC} & \textbf{GFC Lead} \\
\midrule
Institution-type (7 classes) & 7 & $\geq0.70$ & 6 quarters \\
Bank-level (individual) & 30 & 0.805 & 6 quarters \\
G-SIB cross-border & 20 & 0.781 & 6 quarters \\
\bottomrule
\end{tabular}
\end{center}

Super-criticality prevalence: 68.6\% of sample above $\mathcal{C}$.
GFC welfare loss of inaction: 19.57 percentage points
consumption-equivalent. Adversarial robustness: 5 of 12 attack types never
break detection across all configurations.

% ============================================================
\section{CLAIMS}
% ============================================================

\noindent\textit{The following claims define the scope of protection
sought. Where a claim refers to another claim, the dependency is by claim
number. Features of any dependent claim may be combined with any independent
claim or with other dependent claims.}
\vspace{0.8em}

\begin{enumerate}[label=\textbf{\arabic*.}]

% --- Independent Claim 1: DeFi Compliance System ---
\item A computer-implemented transaction compliance system for decentralised
finance networks, the system comprising:
\begin{enumerate}[label=(\alph*)]
\item an interface layer configured to receive transaction requests from
client applications;

\item a compliance orchestration layer comprising:
\begin{enumerate}[label=(\roman*)]
\item a machine learning risk scoring engine configured to compute a
numerical risk score for each transaction request from engineered features
comprising address-level, transaction-level, and graph-level attributes;
\item a graph analytics module configured to construct and analyse
transaction graphs and compute graph-level features;
\item a sanctions screening module configured to check transaction
counterparties against maintained sanctions databases;
\item a state-adaptive transaction flow control module configured to compute
one or more network instability metrics from recent transaction activity and
to select, from a predetermined set of enforcement profiles, an enforcement
configuration comprising:
\begin{enumerate}[label=(\alph*)]
\item a risk threshold profile defining numerical thresholds for a
deterministic compliance decision matrix; and
\item a proof-requirement policy defining, for at least one risk score tier
or transaction class, one or more cryptographic proofs required to be
verified on-chain as a prerequisite to settlement;
\end{enumerate}
and
\item a deterministic policy engine configured to map computed risk scores to
policy actions according to the selected risk threshold profile, and to
enforce the proof-requirement policy by rejecting transactions for which the
required proofs are missing or fail on-chain verification;
\end{enumerate}

\item an offline training layer comprising a composite teacher ensemble
that generates pseudo-labels from unlabelled blockchain transaction data
using a weighted combination of at least two independent teacher models, and
a student model trained on those pseudo-labels; and

\item an infrastructure layer comprising smart contracts deployed on a
blockchain network configured to enforce the determined policy actions before
settlement, including a core contract, an escrow contract, and a
zero-knowledge proof verifier contract;
\end{enumerate}
wherein all compliance decisions are enforced before the corresponding
transaction is settled on the blockchain network.

% --- Independent Claim 2: DeFi Compliance Method ---
\item A computer-implemented method of enforcing compliance in decentralised
finance transactions, the method comprising:
\begin{enumerate}[label=(\alph*)]
\item receiving a transaction request comprising at least a sender address,
a recipient address, and a transaction value;
\item computing a risk score by applying a trained machine learning model
to a feature vector comprising address-level, transaction-level, and
graph-level features;
\item calibrating the risk score using a sigmoid recalibration function
anchored at a percentile of a training score distribution;
\item computing one or more network instability metrics from recent
transaction activity and selecting an enforcement profile from a
predetermined set, the selected profile deterministically defining:
\begin{enumerate}[label=(\roman*)]
\item a risk threshold profile for a deterministic decision matrix; and
\item a proof-requirement policy specifying one or more cryptographic proofs
to be verified on-chain as a prerequisite to settlement for at least one
risk score tier or transaction class;
\end{enumerate}
\item mapping the calibrated risk score to a policy action using the
decision matrix with thresholds from the selected risk threshold profile;
\item enforcing the proof-requirement policy by verifying the specified
proofs on-chain prior to settlement, and applying a non-approval action
when required proofs are absent or fail verification;
\item executing a corresponding smart contract function on a blockchain
network before settlement; and
\item recording the compliance decision and policy action on-chain as an
immutable audit record.
\end{enumerate}

% --- Independent Claim 3: Systemic Risk Controller ---
\item A computer-implemented system for monitoring systemic risk in a
financial network and outputting an enforcement configuration for a
transaction processing network, the system comprising:
\begin{enumerate}[label=(\alph*)]
\item a gravitational interaction engine configured to model a plurality
of financial institutions as particles in a multi-dimensional feature
space interacting through a pairwise potential comprising a bounded
erf-attraction term and a logarithmic repulsion term;

\item a spectral monitoring module configured to compute the spectral radius
of an interaction matrix and to determine whether the network state falls
above or below a critical manifold defined by the condition that the spectral
radius equals one;

\item a blind-spot energy module configured to compute, from reference
distributions fitted on a designated normal period, a complementary energy
function comprising at least an enstrophy anomaly channel, a spectral gap
channel, an alignment anomaly channel, and a temporal novelty channel; and

\item an adaptive policy module configured to compute a state-dependent
friction coefficient from the combined systemic energy and its gradient and
to output an enforcement configuration comprising:
\begin{enumerate}[label=(\roman*)]
\item a risk threshold profile defining numerical thresholds for transaction
approval decisions in a transaction processing network; and
\item a proof-requirement policy specifying, for at least one threshold tier
or transaction class, one or more cryptographic proofs required to be
verified as a prerequisite to settlement.
\end{enumerate}
\end{enumerate}

% --- Dependent Claims ---
\item A system according to claim 1 or claim 3, wherein the pairwise
potential is defined as:
\[
\phi(r) = \gamma\cdot\frac{\sigma\sqrt{\pi}}{2}\,\mathrm{erf}(r/\sigma)
- \gamma\lambda_{\mathrm{rep}}\cdot\log(r+\varepsilon)
\]
where $r$ is the inter-institution distance, $\gamma$ is an interaction
strength, $\sigma$ is an interaction scale, $\lambda_{\mathrm{rep}}$ is a
repulsion strength, and $\varepsilon$ is a regularisation constant.

\item A method according to claim 2 or a system according to claim 3,
wherein the state-dependent friction coefficient is:
\[
\gamma^*(X) \approx \frac{\beta\|\nabla E\|^2}{2\kappa+\beta\|\nabla E\|^2}
\cdot\frac{E(X)}{E(X)+\theta}
\]
where $E(X)$ is the combined systemic energy, $\|\nabla E\|$ is the
gradient norm, $\kappa$ is an economic cost parameter, $\beta$ is a systemic
damage coefficient, and $\theta=\kappa/\beta$.

\item A system or method according to claims 1 to 5, wherein the
enforcement profile is selected from a discrete set of profiles, and
wherein the selection is determined by whether the computed friction
coefficient $\gamma^*(X)$ lies between predetermined lower and upper
thresholds.

\item A system or method according to claims 1 to 6, wherein the
proof-requirement policy comprises, for at least one enforcement profile,
requiring at least one of: a sanctions non-membership proof, a risk range
proof, or a KYC verification proof.

\item A system or method according to claims 1 to 7, wherein the
zero-knowledge proofs are Groth16 zk-SNARK proofs, the sanctions
non-membership proof uses a Merkle tree constructed with the Poseidon hash
function, and each proof has a time-to-live after which it must be
regenerated.

\item A system or method according to claims 1 to 8, wherein the composite
teacher ensemble comprises an autoencoder--gradient boosting teacher, a
rule-based FATF teacher, and a graph-structural teacher, and the composite
score is a weighted sum of their outputs.

\item A method according to claim 2, wherein the sigmoid recalibration
function is:
\[
\hat{p}_{\mathrm{cal}}(x) = \sigma\!\left(k\cdot(s(x)-P_{75})\right)
\]
where $\sigma$ is the logistic sigmoid, $k$ is a steepness parameter, and
$P_{75}$ is the 75th percentile of the training score distribution.

\item A system or method according to claims 1 to 10, wherein selecting
the enforcement profile further deterministically defines at least one of:
a transaction rate limit, a queue delay, an escrow duration, an escrow
release condition window, or a maximum transaction value limit.

\item A system according to claim 3, further comprising a Navier--Stokes
analogy module configured to compute a velocity gradient tensor from
institution position changes, decompose it into strain rate and vorticity
components, compute enstrophy as $\mathcal{E}=\frac{1}{2}\|\omega\|_F^2$,
and evaluate a Beale--Kato--Majda blow-up criterion.

\item A system according to claim 12, wherein the adaptive friction
coefficient is applied as an adaptive viscosity
$\nu=\nu_0(1+\gamma^*(E_{\mathrm{BS}}))$ in the Navier--Stokes analogy.

\item A system according to claim 3, further comprising a Minsky regime
classification module configured to classify the network into hedge,
speculative, and Ponzi regimes based on the cosine alignment between the
gravitational gradient and the blind-spot energy gradient.

\item A system or method according to claims 3 to 14, wherein the
spectral radius is computed using power iteration with finite-difference
Hessian-vector products at $O(N^2 d\log N)$ complexity per time step.

\item A system according to claim 1, further comprising a cross-chain
bridge contract requiring compliance verification on both a source chain
and a destination chain before releasing bridged funds.

\item A system or method according to claims 3 to 15, wherein a welfare
cost of not implementing adaptive friction is computed as:
\[
\%CL = 100\times\!\left(1-e^{-\mathcal{L}_{\mathrm{inaction}}/
(\sigma\theta T)}\right)
\]
and the system operates as an online algorithm with $O(\sqrt{T})$
cumulative regret relative to the best fixed policy in hindsight.

\end{enumerate}

% ============================================================
\section{DRAWINGS}
% ============================================================

\textit{[To be filed separately:
(1) System architecture diagram showing the four compliance layers;
(2) Flow diagram of the compliance decision lifecycle with proof-gating;
(3) Schematic of the zkNAF circuits and on-chain verification flow;
(4) Block diagram of the composite teacher weak supervision pipeline;
(5) Diagram of the GravityEngine $N$-body interaction model;
(6) Plot of the erf-log pairwise potential function;
(7) Diagram of the critical manifold $\mathcal{C}$ in state space with
enforcement profile regions;
(8) Time series of the GFC early warning signal at 2007-Q1;
(9) Navier--Stokes analogy schematic;
(10) Minsky regime classification diagram.]}

% ============================================================
\newpage
\section*{ABSTRACT}
% ============================================================

A state-adaptive compliance enforcement system for decentralised finance
networks enforces regulatory policy before transaction settlement. A
state-adaptive transaction flow control module computes network instability
metrics and selects an enforcement profile; the selected profile
deterministically defines a risk threshold profile for the decision matrix
and a proof-requirement policy gating settlement. A closed-form optimal
friction coefficient derived from a multi-body gravitational model of the
financial network drives profile selection. A machine learning risk scoring
engine trained via composite teacher weak supervision scores transactions;
blockchain smart contracts enforce outcomes before settlement. A
zero-knowledge proof subsystem enables privacy-preserving on-chain
verification of sanctions status, risk range and KYC without revealing
underlying data. The system detects the Global Financial Crisis six
quarters in advance on FDIC data and achieves ROC-AUC of 0.9965 across
2.64 million Ethereum transactions.

\end{document}
